\chapter{为什么造出一组零，就已经快结束了}\label{sec:blocks}
\section{先解决十五元五平均，看看矩阵从哪里来}
设十五个位置已是 $(a^5,b^5,c^5)$，且 $a+b+c=0$。按照下面的表重新分成三组：
\[
K_+=\begin{pmatrix}0&3&2\\1&2&2\\4&0&1\end{pmatrix}.
\]
第一行说“从 $a$ 块拿零份，从 $b$ 块拿三份，从 $c$ 块拿两份”；每行和为五，所以每行可以做一次五平均。每列和也为五，所以恰好用完每个旧块，没有多拿、漏拿或复制位置。

三个新值就是 $K_+(a,b,c)^T/5$。令 $x=a,y=a-b$，则 $b=x-y,c=y-2x$。代入第一行，得到
\[a'=-\frac{x+y}{5};
\]
第二行给 $b'=-x/5$，所以 $y'=a'-b'=-y/5$。这三次真实平均，在两个参数上就是
\[(x,y)\longmapsto-\frac15(x+y,y).
\]
暂时忽略整体比例，它把第一坐标加上第二坐标，这叫\textbf{上剪切}。

另一张表
\[K_-=\begin{pmatrix}2&1&2\\3&0&2\\0&4&1\end{pmatrix}
\]
实现 $-(x-y,y)/5$。两个词接起来，得到 $(x,y)/25$。参数的比例回来了，但真实数值整体缩小了；因此“方向上的逆”完全可以由正向平均实现。

\section{一般的三个等大块}
一般 $p\ge5$ 时写 $p=3m+\epsilon$、$\epsilon=\pm1$。在零和空间上，全一矩阵 $J$ 的作用为零。于是我们可以先设计想要的作用 $I\pm D$，再加 $mJ$，把所有计数变成非负。加上去的部分对零和参数没有影响，却让抽象矩阵成为真的分组。

下面会出现 $\SL_2(\Z)$，它只是所有行列式为一的二阶整数矩阵。$U(k)$ 把 $(x,y)$ 变为 $(x+ky,y)$；$L(k)$ 把它变为 $(x,y+kx)$。欧几里得算法正是在反复做这种整数加减。

$\Gamma_0(3,\Z)$ 是这些矩阵中左下项被三整除的那些。证明中的四陪集表，你可以看成一张四个状态的路线表：读入一个生成元后，记录落在哪个状态，并把差额记成已有的 $U,L$。读完整个矩阵词，就得到可以实际执行的剪切词。
\begin{proposition}[3p个位置的独立解法]\label{prop:threep}
若 $p\ge5$ 为素数，则 $3p$ 个位置满足定理\ref{thm:prime}的判据。
\end{proposition}
\begin{proof}
先将位置分为三个 $p$ 组。若其和模 $3$ 全相等，交换不同组中一对不同余值，组和变为 $s+\delta,s-\delta,s$，其中 $\delta\ne0$。这样的跨组数对必存在，否则原配置模 $3$ 全相等。分别平均后得到 $(a^p,b^p,c^p)$，$a+b+c=0$，且模 $3$ 不全相等。

写 $p=3m+\epsilon$，$\epsilon=\pm1$、$m\ge2$。令 $J$ 为三阶全一矩阵，
\[
D=\begin{pmatrix}1&-1&0\\1&-1&0\\-2&2&0\end{pmatrix},\qquad
K_\pm=mJ+\epsilon(I\pm D).
\]
两表非负且行列和为 $p$。每个输出组按一行从三个旧块取数，三次平均实现 $K_\pm/p$。在零和空间上，它们等于 $\epsilon(I\pm D)/p$。由于 $D^2=0$，有
\[
(I+D)(I-D)=I-D^2=I.
\]
实际三步平均仍会带上 $1/p$ 的共同尺度，但两作用在共同尺度意义下互为逆。
因此两个方向的整数次剪切都有正向平均序列。
改用 $(x,y)=(a,a-b)$，它们分别为 $U(1),U(-1)$。块角色交换给
\[
\Sigma=\begin{pmatrix}1&-1\\0&-1\end{pmatrix},\qquad
P=\begin{pmatrix}-2&1\\-3&2\end{pmatrix},\qquad
L(-3)=U(-1)P\Sigma.
\]
因此 $U(\Z),L(3\Z)$ 可实现。


这里 $\Gamma_0(3,\Z)$ 表示行列式为 $1$、左下项被 $3$ 整除的整数矩阵群。
为什么恰好出现 $3$？因为我们正在控制三个等大块，
其模 $3$ 的差分就是剩余的共同余数障碍。
这些矩阵射影生成 $\Gamma_0(3,\Z)$。为验证，取 $S=\left(\begin{smallmatrix}0&-1\\1&0\end{smallmatrix}\right)$、$T=U(1)$，右陪集代表为 $I,ST^{-1},S,ST$，由模 $3$ 底行的四个射影类区分。
陪集表的含义是：从代表 $r$ 出发，右乘一个生成元 $g$，
再乘新陪集代表 $r'$ 的逆，便回到子群。
任意群词可以逐字这样改写为这些回路的乘积，
因此只要表中的回路都由我们已有的矩阵生成，整个子群就由它们生成。
对生成元 $S,T,T^{-1}$，Schreier 回路 $rg(r')^{-1}$ 如下：
\[
\begin{array}{c|ccc}
r&S&T&T^{-1}\\\hline
I&I&T&T^{-1}\\
ST^{-1}&-T&I&L(3)\\
S&-I&I&I\\
ST&T^{-1}&L(-3)&I
\end{array}
\]
欧几里得消元证明 $S,T$ 生成 $\SL_2(\Z)$，上述表便给所需子群；中央符号不影响式\eqref{eq:return}。

本原三块的 $G=\gcd(y,3)$。合法性给 $3\nmid y$，可在 $\alpha x+\beta y=1$ 的裴蜀解中令 $3\mid\alpha$。矩阵
\[
\begin{pmatrix}y&-x\\\alpha&\beta\end{pmatrix}\in\Gamma_0(3,\Z)
\]
将 $(x,y)$ 送到 $(0,1)$，故得到 $(0^p,s^p,(-s)^p)$。
到了相反值配置，令 $h=(p-1)/2$。第一次选 $h$ 份 $s$、$h$ 份 $-s$ 和一个零，
得到 $p$ 个新零；第二次作同样的选择以后，每个符号只剩一份。
最后一组用这两份相反值和 $p-2$ 个零即可。
具体地，连续两次平均 $(s^h,(-s)^h,0)$，最后平均 $(s,-s,0^{p-2})$。
\end{proof}

\section{一个零给我们留了什么余地}
在 $2p+1$ 个零和位置中，如果有一个零，把它留着，其他 $2p$ 个位置正好分两组平均。两个组和相反，所以得到 $(a^p,(-a)^p,0)$。

例如 $p=5$，先拿两份 $a$、两份 $-a$、一个零平均，产生五个零；再做一次，只剩一对 $a,-a$；最后拿这对和三个零。只需三次就消去两个五重块。
\begin{lemma}[含零的2p+1元配置]\label{lem:zero-base}
对奇数 $p$，$2p+1$ 元零和配置若含一个零，则至多五次 $p$ 平均归零。
\end{lemma}
\begin{proof}
保留零，将其余位置分成两个 $p$ 组并平均，得到 $(a^p,(-a)^p,0)$。用命题\ref{prop:threep}末尾所列的三次相反值操作完成。
\end{proof}

\section{为什么可以在较小子系统里收尾}
一个零和 $p$ 元组平均以后，至少出现 $p$ 个零，所以非零位置最多有 $n-p$ 个。如果我们选一个子系统，\emph{包含全部非零位置，再多带一个零}，它的总和仍然是零，目标与整个配置一致。

含零还带来一个算术好处：所有坐标差的最大公因数等于坐标本身的最大公因数。因为差里面就包含 $x_i-0=x_i$。本原化后，这个子系统的 $G$ 必为一。

因此，只要这个子系统的大小属于已经证明的范围，就能在里面完成，外面的零保持不动。下面按 $n$ 的范围选择 $2p+1,3p,4p$ 等已解决的大小。
\begin{lemma}[从一个零和组到全部归零]\label{lem:zero-tail}
若 $p\ge5$，$n\ge2p+1$ 的零和配置含一个零和 $p$ 元组，则可归零。对任意奇数 $p$，此结论在 $2p+1\le n\le3p$ 同样成立。
\end{lemma}
\begin{proof}
先平均该组，产生 $p$ 个零。若某个已解决的维数 $m\le n$ 满足非零支持数小于 $m$，就选一个 $m$ 元子系统包含全部非零位置及一个零。它自身零和，本原化后的差分公因子为 $1$，可在该子系统归零。
对 $n\le3p$，使用引理\ref{lem:zero-base}中的 $m=2p+1$。对 $3p<n<4p$，使用命题\ref{prop:threep}。对 $4p\le n<5p$，使用引理\ref{lem:integer-range}在 $m=4p$ 的结论，因为 $d_p(4p)=1$。对 $n\ge5p$，同一引理的最终上界已适用。
\end{proof}

后面每次说“送到终端”，具体目标都是先让某个 $p$ 元组的和恰好为零。其余操作由本讲接手。这样我们不必要求一个可逆矩阵把非零列直接变成零；它也确实做不到。
